\chapter{粘性流体力学基础}
实际流体都是有粘性的，只有当粘性应力很小时，才可忽略粘性影响而采用无粘 流假设的"理想流体"模型.本章介绍粘性流体运动的主要性质、支配其运动的基本方 程以及若干典型流动的求解方法.

\section*{8.1  粘性流体的本构方程}\addcontentsline{toc}{section}{8.1 粘性流体的本构方程}

在第 3 章讨论流体运动基本方程的封闭性时曾指出 z 流体运动的 3 个守恒定律 只有 5 个独立方程，而未知数有 12 个，需要补充状态方程(1个)以及应力张量与流 体运动间关系的本构方程 (6 个)以后，基本方程才能封闭. 导出流体本构方程有两种方法 z ${\rm 1}$基于分子运动的统计力学方法;${\rm 2}$理性力学 和实验相结合的方法.力学和工程中常用第二种方法.各种各样流体的物性千差万 别，就流体应力状态和流体运动的关系来说，可以分为两大类，具有记忆性的流体和 非记忆性流体.流体质点的应力状态和该点的变形及应力历史有关的称作记忆性流 体，例如稠油和高分子溶液.非记忆性流体的质点应力状态只和当时该点的流体变形 状态有关，或者说，流体质点的应力状态对于流体变形是瞬时响应的.下面只考虑非 记忆性流体的本构方程. 1.建立本构关系的基本原理 本构方程属于物性方程，应当具有普适性.根据理性力学原理，建立本构方程应

当符合以下基本原则. (1)可表性原则 应力和变形率都是张量，因此本构方程应是张量方程，它应当遵循张量运算规 则.例如质点的应力 T'j 是二阶对称张址，微团的变形率丘，也是二阶对称张量.关系 式 T'J= $\alpha$ S'J (a 是标量〉和坐标系无关，它在任何坐标系巾都是二阶对称张量之间的 关系式，它符合可表性原则.但是速度梯度的反对称张量。"和质点应力 T 'J 之间的简 单线性关系是不成立的，即 T'J 手i:.\{3n 'J (${\rm ß}$ 是标量) ，闲为一个对称张量不可能和 一 个反 对称张撞在所有坐标系中相等，所以它不符合可表性原则.又如， 一 个对称张量 T'J 和 一个向量 Uj 之间的关系 T ij = $\alpha$ ijk Uk 可以成立的条件是 :$\alpha$ ，jk 必须是 三 阶张量，并且 $\alpha$ ，片关于下标 i .j 对称.总之，粘性流体的本构方程必须符合张璋的运算规则. (2) 客观性原则 本构方程是物性关系，因此它应当和参照系无关.例如.流体的本构方程和它的 状态方程 一 样，在地球和月球上应是相同的.具体来说，若有两个观察者.一个在图定 的参照系中，而另 一 个在相对运动的参照系中.这两个观察者得到的本构方程应当相 同.由于参照系的运动可用坐标架的刚体运动来描述，它有 三 个平移和 三 个转动自由 度.因此客观性原则可表述为，本构方程的函数表达式应不依赖于参照系的刚性运 动.应用客观性原则，可以导出本构方程应有的基本形式. 设参照系统刚性运动方程可以表示为 x' = x" (t) + Q(t) ${\rm •}$ x (8. 1 a) 或

\begin{equation*}.2< = X"i (t) + Qij (t)Xj\tag{8.1b}\end{equation*}

式中 x' 是运动坐标架中任意点的位置向量 .x 是固定坐标架对应点的位置向量， Xo (t) 是两坐标架间的相对位移向量，它只是时间的函数 . Q(/) ( 或 Q'J) 是坐标架的转 动张量，转动张量的具体表达式可由任意刚性坐标系的变换式求出.设有两个直角坐 标系，当坐标系 eJ 相对于 e ， 作刚性转动时，这两个坐标基之间有以下关系:

\[, , ei = \alpha );e j' e, = \alpha  'JeJ\]

直角坐标基变换系数的是对应坐标轴间的方向余弦(见表 8. 1 和图 8. 1) ，并有以下 性质 z $\alpha$ liQI\} - 此， $\alpha$""$\alpha$ 1'" = $\delta$'J 褒 8.1 坐标轴阎方向余弦 e l e 2 e :i

\begin{gather*}
e , l \alpha 11 \alpha 1 2 \alpha 13\\
e , 2 \alpha  2 1 U :!2 \alpha 2 3
\end{gather*}

e , z $\alpha$ 31 $\alpha$ :i2 $\alpha$旧

两坐标系中位置向量的变换有以下的等式 z 3.:; =酌iX j' X 1 = $\alpha$ ijX J 对照式 (8.1b). 转动张量 Q(t) 等于当时的坐标基向量 e; 的 方向余弦张量，即 Q$\upsilon$=$\alpha$。或 Q(t) = $\alpha$ i， ejej Q (t) 或 Q。有以下的性质: Q\{, Q/j = 8i,. Q刷 Q"" =${\rm O}$Ij或 Q.QT=1 z 图 8. 1 坐标系的变换 根据张量性质，任意一个二阶张量扎在不同坐标系中应有以下等式 z $\rho$;=$\alpha$d$\alpha$ 川户 /111 因此在坐标架作刚性转动时，二阶张量丸，应满足以下关系 z P:j = Q\textasciitilde{} 户阳 Qjm 利用以上的关系式.根据客观性原则可以导出本构方程的基本形式.假定在固定 参照系 e ， 巾有一个张量函数关系式(已符合可表性原则) :

\[r" = fij (.p .B/ .A...)\]

h 表示张量函数 J 表示标量、 B/ 表示向量、 Anm 是二阶张量.在运动参照系 e: 中上式 应写成

\[r:j = f:, (.p' .B~ .A'....\cdots )\]

根据前面导出的关系式，在坐标架刚性转动时，以上两式中各量应满足以下等式 z r:j = Q\textasciitilde{} (t )r/m Q 呵( t) ${\rm •}$ f:j = Q ,/ (t) f阳 Q町 (t) ${\rm •}$ 骨， = .p. B: = Qij (t) 矶 ${\rm •}$ A:, = QiI (t) A 阳 Qmj (t) 客观性原理要求，不论参照系作何运动，物理量 T;j 和札 B 、 A 之间的本构函数关 系应始终保持不变. 然而并不是所有的物理量都能满足客观性条件，例如，速度就不满足客观性条 件.对位置向量式 (8. 1)求时间的导数，可以获得运动参照系中速度表达式 z U: 二二:= .i\textasciitilde{}， (t) +Q必 .i k +Q$\mu$ 矗 = U Oi(t) + Q必 Uk + Q${\rm u}$. Xk (8.2) 上式右端有 UOi (/)丰nQ"x， .它们和坐标架的运动有关，它不符合 U;=Qji (t) U尸因此， 速度不能作为自变量单独出现在张量型本构方程中.另外还可以证明，速度梯度张量 aU ，I J .:r j 也不是客观量，但是对称张量$\theta$ ui/aXj 十 cJ U;/ ${\rm e}$J.:r j 是客观量.利用刚性转动坐 标架中速度公式 (8.2) 计算向量梯度，因为 .i $\tilde{i}$ (t)和 Qij 只和时间 t 有关，于是可得 丝;=Q\& 冉 +Qd44

\[<l.:r, f'.:r, f'.:rj\]

cJ x/ a $\tilde{a}$ 注意到梯度算符工-，=一丁一 =Q 一，因而速度梯度张量r1.1;'- iJ.:r;' <1.1'/ - ""j/ r1 :l'/ 在=叫苦 )QJ/ +Q公告 =QdZfQJI+QAJh

上式右端有 QikQj是，它不符合张量关系式 A\textasciitilde{}=QJ (t)A/..Q ", j (t) ， 因此二阶张量 ${\rm e}$J U'; ${\rm e}$J Xj 也不是客观量.同理，计算二阶张量 ${\rm e}$J Ui/ ${\rm e}$J Xj 的转置 ${\rm e}$J U;i${\rm e}$J Xi' 可得 JU'i '" 'au/ '" ,"; JXk '" JU 了千 = Q j/ ':U/Qik +Qj' 了\textasciitilde{}=Q品 $\tilde{j}$/+QjkQikCJX i - cJ X矗 cJX i cJX, 将上两式相加，得 JU: I au\textasciitilde{} '" \{${\rm e}$JU/ I ${\rm e}$J U川1-+一$\rightarrow$ =Qikl 一-'/ +\textasciitilde{}一!!. JQ,/ +Qjk Qik +Q必 Q必JX'i ' Jxj 必飞 ${\rm e}$J Xk aX/ I\textasciitilde{}" ' \textasciitilde{}，. $\tilde{W}$ ' $\tilde{J}$< 由于转动张量是对称正交张量，即有 QjkQik = 鸟， 将上式对时间 t 求导数，得 Q，kQik +QjkQ. =0 ， 于是前式最后两项之和等于零.也就 是说对称张量au; /cJ X: + $\theta$U:/ai 满足客观量条件:

\[eJ U~ , eJ U~ - l eJU , , aU.  \]

74 十 7千 =Q必 l 了\textasciitilde{}+-...-二 JQ;/ cJ Xi cJ Xj 飞 dX矗 dX/ J 用向量算符表示，对称张量aU'; $\theta$ Xj +aUj !aXi ， 也可写作 (VU)+(VU)T. 在流体 运动中，已经导出速度梯度的对称张量之半是流体运动的变形率鸟，即 5ij = \textasciitilde{} (但\textasciitilde{} + \textasciitilde{}$\tilde{J}$) 2 飞 JXj , aXi I 所以变形率张量也是客观量，因此如下应力张量和变形率张量之间的本构关系式符 合可表性和客观性原则(为了避免和温度符号 T71昆淆.应力张量暂时用 P 表示) : Pij = Iij (25阳) 对于不同流体，它们有不同的函数关系式 fij. 总之，根据客观性原则，可以导出本构 关系式中变量的形式.但是客观性原则不能导出本构方程的具体函数关系式，具体的 函数形式依赖于对物性的分析和实验.下面较详细地介绍牛顿流体的本构方程. 2. 牛顿型流体的本构关系 定义 8.1 粘性应力张量 P 和变形率张量 S 间具有线性各向同性函数关系的流 体称为牛顿型流体. 例如:常温常压下的空气、水、酒精和稀油等都属于牛顿型流体，牛顿流体是一 种简单的非记忆性流体.下面来导出牛顿型流体的本构关系式. (1)粘性流体的应力分解 z 根据张量运算的规则，流体中任一点的应力张量 (2 阶张量)必可分解为一个各向同性张量和另一张量之和 z

\begin{equation*}P =- III +r\tag{8.3a}\end{equation*}

或用分量形式表示为 Pij = 一 [J(J ij + Tij (8.3b) 或将它写成矩阵形式

\begin{gather*}
P 11 P I2 P 13 1 (-11 0 0 1 r<1I <1 2 <13\\
P21 P22 P23 I = I 0 - 11 0 1+ 1<21 <22 <23\\
P 31 P n P 33 J 1 0 0 - 11 J 1 <31 <32 < :l3
\end{gather*}

式中 E 是应力的各向同性部分，为标量函数.当流体静止时，流体内部不存在粘性应 力，这时 E 应等于静止流体的压强 $\rho$ ，也就是热力学压强;流体运动时 E 不等于热力 学压啤 $\rho$ . Tij 称作偏应力张量，当运动停止时 ， r;J 应趋于零.因 P i， 是二阶对称张量， . 118ij 也是二阶对称张量，所以由式 (8.3) 可得出结论 : T，j 也是二阶对称张量.下面先导 剧本构关系中应力的各向同性部分. (2) 本构关系中各向同性应力的线性关系式 流体运动状态下各向同性应力 H 与热力学压强 $\rho$ 不相等，但当流体运动停止 后 .11 趋于静止流体的热力学压强，因此可以把 E 分成两部分 z 118ij = ($\pi$ - pMi, = 与流体运动有关部分+热力学压强 其中 $\rho$ 是热力学压强，它是流体的状态参数，和流体的变形率没有直接关系.根据线 性假定 $\pi$ (5 ij ) 应是流体运动变形率张量 5ij 的线性函数.因为 $\pi$ 是标量，而岛的线性 标量函数只有它的迹 5 kk ( 即 5 ij 张量的对角线项之和 ) .故必有 $\pi$ =${\rm A}$5kk ${\rm •}$ 从而牛顿型 流体的各向同性应力张量 H 应为 118i, = ($\rho$ 十 ${\rm A}$5 kk 用。 (3) 偏应力张量和变形率张量间具有线性各向同性关系 偏应力张量是二阶对称张量，上面根据建立本构关系的原理已证明，非记忆性 流体的本构关系应有<" = !ij (5ij ) 的形式，各向同性线性画数 !ij 的唯一形式是 z

\[<i, = 2\mu 5"\]

式中 $\mu$ 是标量，称为粘性系数，将由实验确定. (4) 牛顿型流体的本构关系 将各向同性应力和偏应力张量相加，得牛顿型流体的本构关系式如下 z Pij = II${\rm o}$ ij + <ij = (一 $\rho$ + ${\rm A}$5kk Mij + 2$\rho$" 为了和常用粘度的定义 一 致，引用符号 j=A 十 2$\mu$/3 ，则 A=j 一年/3 ，于是上式可改 写为 mo= - H"-3向8'1 +向8 ij 牛顿流体本构关系的最终形式为 Pij = [一 $\rho$+(j-3$\mu$) 叫8 ij 十队 (8.4) 式 (8.4) 也称为 "r 义牛顿定律"式中 $\mu$ 称为第一粘性系数，或动力粘性系数，简称粘 性系数 zj 称为第二粘性系数.我们知道 5kk =511 +522 +533 =JUi/Jx , =V ${\rm •}$ U ， 它是

流场的散度，也是质点的体积膨胀率.因此本构关系式 (8.4) 表明:牛顿型流体的质 点应力状态由以下 =三 部分组成. (- p8ij : 热力学压强; ${\rm 2}$ ($\mu$'-34Moz 由体膨胀率(或压缩)引起的各向同性粘性应力; ${\rm 3}$ 2$\mu$ 吼:由运动流体变形率引起的粘性应力.称为偏应力张量. ${\rm •}$ (5) 粘性系数 $\mu$ 的确定 根据连续介质原理和线性假定，已经导出牛顿型流体本构关系的一般形式，但是 其中有两个物性系数 $\mu$ 和 $\mu$' 是待定的，宏观的流体力学用实验方法测定这些系数. 通常确定粘性系数的方法属于流变学的范畴，测试的基本原理是，构造 一 种简单流 动，使它只有一个剪应力和剪应变率.然后分别测定剪应力和剪应变率来得到粘度 系数. 例如:用很大的两块平行平板，下板固定，上板以平行于下板作直线匀速运功， 这样就可以构造一种特别简单的剪切流动(见图 8.2): U) (X2) ,U2 =u\textasciitilde{} =0. X14 u 、 l f) 图8.2 简单的平行剪切流动 此时流场的剪应变率只有: S)\textasciitilde{} = (du) / dX2 ) /2. 测定上版面的剪应力.就可以代 人本构关系式来确定粘性系数.理想的做法是设计平行平板间的距离 d 很小，因此 平板间流体的速度分布几乎是直线，即 u) =U.l' 2 /d( 参见图 8.2) .于是 S)2=U/2d. 同 时测定上饭面的拉力 T 和上极的面积 A. 则上版面的剪应力 <12 = T/ A. 最后得粘度 系数

\[\mu =2385)\]

牛顿曾在流体作简单直线层流运动情况下，假定剪应力与剪应变率间关系 式 (8. 日，因此这种流体称作为牛顿流体.在非均匀剪切流动1f t .式 (8.5) 可推广为

\[u-Y EG\]

丁。$\mu$ 一-r (8.6) 式 (8. 5) 或式 (8.6) 称作牛顿粘性公式， $\mu$ 称为流体的动力粘性系数，简称为粘度， $\mu$ 的单位是帕$\cdot $秒 (N. s/m 2 ). 丁.程中还常用动力粘性系数除以密度(即 $\mu$ /p) 作为流 体粘度的度量，称为运动粘性系数，用 $\nu$ 表示，单位是时 /s.

(6)$\mu$' 的物理意义 现在来考察 $\mu$' 的物理意义.任取一固定点 M ，考察以 M 为中心，以 a 为半径的无 限小球面 A 上的法向应力的平均值 P m ， 可以证明 Pm=Pii/ 3 . 根据定义 z Pm=lim 」$\tau$ dJ) P. ${\rm •}$ ndA = lim \textasciitilde{}$\tau$ <[${\rm þ}$ n ${\rm •}$ P ${\rm •}$ ndA .-0 4$\pi$a- JJ " o 4$\pi$a- JJ

\[=lim~\tau  <[\theta  n in iP ii dA\]

u  o 4$\pi$a- JJ A 因为半径 a 无限小，可以将 P ij 取作 M 点的值，因而移到积分号外.同时球面的法向 量可表示为 ni =.r ;la ， 于是 Pm=lim 」$\tau$ ${\rm Æ}$.${\rm Þ}$ n ， njP II dA = lim 乒.${\rm ÆÞ}$ 马 jdA "() 4$\pi$a- JJ "() 4$\pi$a- JJ a 利用高斯公式不难证明: 于是，有 Pm= lim 乒与吵0 王L与F "() 4$\pi$a矿- JJ$\alpha$"户$\rightarrow$() 4$\pi$" 叼」叩\#dz鸟J P;; 4 71" 月 3 =lim 一」$\tau$. ，$\tilde{w}$. 8;i = : P;;

\begin{gather*}
"\rightarrow  04\pi a' 3 \upsilon 3 - ..\\
Pm=\div (Pll+P22+Pm) (8. 7)
\end{gather*}

这表明 M 点处正应力的平均值等于三个相互垂直坐标轴方向正应力的平均值.根据 定义，它是一个不随坐标系变化的不变量，将式 (8. 的代入式 (8.7) 可得 Pm = 一 $\rho$ 十，/ v. u 由上式可看出: ${\rm 1}$对不可压缩流体: V ${\rm •}$ U=O ， 故 Pm=-P ， 此时本构关系中含 $\mu$' 的项为零，就 是说，不可压缩流体一点法向应力的平均值等于热力学压强. ${\rm 2}$对于可压缩流体:在运动过程中，流体微团的体积发生变化，它将引起平均 压强 Pm 值发生变化，其变化量为 $\mu$'v ${\rm •}$ U ， 故 $\mu$' 称为"容积粘性系数"或"第二粘性系 数"，由此可见， j 反映由体积变化引起流体偏离热力学压强的粘性应力.但实际情况 是:除高温和高频声波等极端情况外，一般的气体运动中可近似认为 j=0. 假设 j=0( 常称斯托克斯假设) .得到牛顿型流体本构方程的简化形式为 Pij = 一到ij +与 (so-t$\omega$$\upsilon$) 在直角坐标系中，式 (8. 的的分量形式可写为 $\iota$=-p-3川 ${\rm •}$ U+2$\mu$ 去 $\rho$"=-p-3川 ${\rm •}$ U+ 与 (8.8)

\begin{gather*}
h-h\\
\mu  nL + U V \mu 
\end{gather*}

2 一3

\[ba --9\cdot \]

、、${\rm ‘}$，，，，，w-z a 之 d

\[+ u-z qo--10 , , \]

EEEE 飞、 $\mu$ 一-P 一-P PF=Pzy=42+ 窍) PZ.y = Pyr =$\mu$( 艺+去) 对于不可压缩流体 V ${\rm •}$ U=Q ， 式 (8.8) 可简化为

\begin{equation*}Pij =- "\theta 8ij +2\mu Sij\tag{8.9a}\end{equation*}

或 P =:- pl + 2川 3. 非牛顿流体的近似本构方程 定义 8.2 不满足牛顿流体本构关系的流体称为非牛顿流体. 悬浮液、聚合物溶液或原油、水泥浆、血液等都是非牛顿流体.非牛顿流体的本构 关系较复杂，尤其是具有记忆特性的粘弹性液体更复杂.下面介绍较简单的非记忆特 性的拟塑性流体(即无运动历史效应的非牛顿型流体)的本构关系. (1)幕律流体 幕律流体和牛顿型流体的主要区别是 z 它的偏应力张量和变形率张量间不是线 性关系，而有以下本构关系式: (8.9b) r 1 寸 2 '[' = /(2S) = kl 言 tr(2S)" I (2S) 式中 tr 表示张量的迹，例如: trA=A 们幕律流体的本构关系还可表示为 '['=协 $\mu$ = k[ $\tilde{tr}$(2S)2 ]早 (8. 1 Q) 式中 $\mu$ 称为表现粘度 ， k 、 n 为常数，由实验测定. S = [vu + (VU)TJ/2 , tr( 2S)2 是 (2 的 2 的迹，等于 (2S)2 张量对角线上三项之和.大多数泥浆、水泥浆、非稠油等可近 似为幕律型非牛顿流体. 幕律流体的表现粘度也可由简单剪切流动 (u=u(y)) 测定.这时剪应力 r 和简 单切变率${\rm y}$ (=du/dy) 之间有以下关系 z

\[r=\mu  (j) i\]

$\mu$( 苦)称为表现粘度，它具有以下的幕画数形式=

\[\mu  (j) = k y<' 1) /2\]

图 8.3 表示不同事指数下剪应力和剪切变形率的关系.指数 n=l 就是牛顿流 体，这时剪应力和剪应变率是正比关系川 >1 的事律流体称为剪切稠化流体，此时，

\section*{8.2  牛顿型流体的运动方程 s 纳维·斯托克斯方程}\addcontentsline{toc}{section}{8.2 牛顿型流体的运动方程 s 纳维·斯托克斯方程}

当剪切率增大时，罪律流体的表现粘度大于牛顿流体的粘度 z 与此相反 ， n<l 的事律 流体称剪切稀化流体. (2) 宾祝流体 宾汉流体的本构关系如图 8.3 上曲线 D 所示. 当流体的剪应力小于某一值时，它呈现固体特性， 没有运动;当剪应力大于某值口以后，它像牛顿流 体一样流动.它的本构方程可近似为 (8.11)

\[n=1\]

剪切率?当 1.1> 口， .-.o =$\mu$ j; 当 1 ${\rm •}$ I\textasciitilde{}.o . ${\rm y}$ = 0 图 8.3 粘性流体的流变特性曲线 ?为剪切变形率.

在第 3 章中，从质量守恒、动量守恒、能量守恒原理由发，得到了流体力学微分形 式的基本方程组 z 3+V$\cdot $( 卢) = 0 (8 山 田J \textasciitilde{}， 1 言\textasciitilde{} = f + $\tilde{V}$ ${\rm •}$ P (8.13) $\upsilon$I P

\[1 ~ , ~..- 1 57(e+UZ/2)=f-u+-V\cdot (P-U)+-v-UVT)+q(8.14 >\]

p p 对于牛顿型流体运动，应用 8.1 节建立的牛顿流体本构方程 (8. 的，将它代入运动方 程和能量方程，就可得到牛顿型流体运动的封闭方程组. 1.纳维-斯托克斯( Navicr-Stokes) 方程 将牛顿流体的本构方程 (8. 的代入式 (8. 13). 得牛顿流体的运动方程 军 = f-\_!\_VP-\_!\_V ( $\iota$V.U)+$\tilde{V}$.(2$\mu$S) Ul P P 飞，$\tilde{I}$ 式中左边是流体质点的加速度，右边各项的物理意义是: f 为作用在微闭上单位质量流体的质量力; \_\_!\_V$\rho$ 为作用在微团上单位质量流体的压强合力:p -tv(3$\mu$V ${\rm •}$ U) 为作用在微团上单位质量流体的粘性体积膨胀力 z tv$\cdot $(2川)为作用在微团上单位质量流体的粘性偏应力张量之合力 (8. 15)

288 第 8:章 粘性流体力学基础 在直角坐标系中牛顿流体的连续方程和运动方程可写作 z (J p \_1 $\theta$( 卢) - " 一一一  t (J Xj 毕 =f-iELii\{ 主$\mu$ 旦)+.l ';-I J1(坐+旦门 (8. 17) Dt J' P dXi P (J Xi \textbackslash\{\} 3 t' (J Xj J $\rho$ aXj L..... \textbackslash\{\} dXj , iJXi J J 2. 不可压缩牛顿型流体的连续方程和运动方程 对于不可压缩流体，它的体积膨胀率等于零，即 $\theta$ Ud iJ Xj =0 ，常粘性系数牛顿流 体的控制方程可简化为 它的向量形式方程为 匹\textasciitilde{} = fi 一土豆主+旦土(坐\}P iJXi ' P (J xJ 飞 ch' ， 1 I iJUj \_ " iJ Xj 要 =f-ivp 十 "${\rm o}$.U UC P

\[V. U = 0\]

式中 $\nu$=$\mu$ /p 为运动粘性系数，通常它是常量. 3. 不可压缩牛顿流体运动的能量方程 在直角坐标系中将应力做功项展开，得 p. U = P;jeiej ${\rm •}$ U.e. = PijU.$\tilde{j}$.e; = $\rho$ ijUje; V. (p. U) = 乌兰 . (PiiUie i ) OX. =\textasciitilde{}.i 主 (PiiU i ) = 主 (PiiU ，) cJX. OXi 将上式代人式 (8.14) ，得到直角坐标系中粘性流体的能量方程 (8. 18a) (8. 18b) (8. 19a) (8.19b) D I L U2 \textbackslash\{\} $\tilde{Tr}$， 1 ${\rm e}$J, $\tilde{Tr}$, , iJ I $\theta$ T\textbackslash\{\} ::.1 e + \textasciitilde{}...， 1 = fV霄+一-一 (P;;U;) + 一:\_ 1 ).一一 1+ q (8.20) Dt 飞 2 I J ,-, , P dXi ,- 'J-J' , (J Xi \textbackslash\{\}.. iJ Xi I ' '1 式中各项的意义是:

\[D I , U2  \]

-le+-l 为单位质量流体总能量的增长率;Dt \textbackslash\{\} -, 2 I fi U 为质量力做功; 士未叩j) 为表面力做功 3

J: i (A \textasciitilde{}\textasciitilde{})为热传导输入能量 z d 为其他能源(如化学反应热等) . 能量方程还可以写成其他形式.以 U 点乘运动方程 (8. 19b) 得 再用式 (8. 20) 减上式得 因 Pij 是对称张量，故 DUi rYT' 1 ap Ui $\tau$\textasciitilde{} = fVi + 一 $\tau$'.....!l.. Uj

\[Ul P <1.Ti\]

De 1 $\tilde{aUi}$ , 1 d 1 , ${\rm e}$J T 飞一= $\tilde{p}$ 一:::.L+一一:\_ (A 一-=-1+ ci Dl $\rho$') dXi ' P JXi 飞 JXi J ' ... ${\rm e}$JU i T< r 1 \{aui I ${\rm e}$JUi \textbackslash\{\} l P ii 了:1.= P ii I$\tau$( \textasciitilde{}\textasciitilde{}-) + :, -, 1 I ${\rm •}$ dXi . L L.飞 <1 X ， cl .Tj I J

\[= (- POij +2\mu Sij )Sij = 2\mu SijS \eta  -pSii\]

对于不可压缩流体.上式最后一项等于零，并令 $\phi$=2$\nu$S$\upsilon$ 乱，将它代入能量方 程后，得 De 1 ${\rm e}$J l ，i. 丁飞 =三 =\textasciitilde{}-;;-($\lambda$$\div$二二 l 十 ci+ $\phi$ (8.21a) Ut p d .T i 飞 d.Ti I 式中 $\phi$=2 凶ij Sij >0 称为耗散函数，它总是正值，是流体内部粘性耗散的能量. 式 (8. 汩的说明.单位质量流体内能的增长率等于热传导输入的能量、生成热和粘性 耗散能量之和.根据热力学惰的定义， TdS=de+$\rho$dO/p) ，对于不可压缩流体， dp= 0 ，因此 TdS=de ，能量方程 (8.21a) 又可写成 DS ., 1 (J \{, a T \textbackslash\{\} 一一 =ci+ 一一:\_ (;.号$\tilde{I}$+$\phi$ (8.21b) Dl $\tau$$\rho$ cJ.T i 飞 <${\rm I}$ Xi I 该式说明:流体质点恼值的增长率等于生成热、热传导输入的能量和粘性耗散的能 量.式 (8.21a) 和式 (8.21b) 说明:粘性耗散使不可压缩流体质点的内能和摘值增加. 方程式 (8. 16) 、式 (8.20) 和式 (8. 21 )构成了牛顿型流体的封闭方程组. 4. 定解条件 有了粘性流体运动的封闭方程组，求解流动问题还需要初始和边界条件.粘性流 体流动的定解条件如下. (1)非定常流动的初始场 求解一般非定常流动问题，需要给出流场的起始状态，也就是

\begin{equation*}1 = O:U = U(x ,  O) ,  \rho =\rho  (x ,  O) ,  T = T(x ,O)\tag{8.22}\end{equation*}

(2) 边界条件 ${\rm 1}$固壁粘附条件，或元滑移条件 在固体壁面上，流体粘在壁面上而无滑移，所以流体速度与当地壁面速度相同 z

\[U = Uh\]

在绕固定物体的定常流动中，壁面速度肌 =0. 从而无滑移条件为 z 壁面上

\[U = o\]

${\rm 2}$固壁温度的热平衡条件 流体质点温度应与当地壁面温度相同，即

\[T w = T1,\]

${\rm 3}$绝热固壁条件 绝热固壁的热传导等于零，即 (某人 = 。 ${\rm 4}$流体分界面上的运动学、动力学和热力学条件 (8. 23a) (8.23b) (8.24a) (8.24b) 在互不侵入的两种流体分界面上，如不计表面张力，界面两侧任 一 点流体的速 度、应力和温度(分别用下标"+"、"一"表示)应相等，即有

\[U1 = U-. P 1 = P. T , = T\]

5. 粘性流体运动的基本特性 (1)粘性流体运动的有旋性 (8.25) 可以证明满足不可压缩理想流体无旋流动的解也满足纳维 - 斯托克斯方程.因 为无旋流动的速度场满足欧拉方程 z DU 1 古- -一 -=-Vp ， V.U=Q l .ll P 而且它的速度场是有势的 U=V$\phi$ ，速度势满足拉普拉斯程 ，1:::. $\phi$ = 0. 将该解代入 式 (8. 19a) 和式 (8. 19b) ，由于 I:::.U= I:::. (V$\phi$) =V (1:::. $\phi$)=0. 因此，纳维 - 斯托克斯方程 中的粘性作用力项等于零，也就是说，无旋流动的速度势解可以既满足欧拉方程，又 满足纳维-斯托克斯方程.但是无旋流动的解在壁面上只满足不可穿透条件:

\[J\phi A\]

()n 而切向速度 U， = $\theta$$\phi$/ () s 手 0 ，也就是说无旋速度场的解不能满足固壁无滑移条件 式 (8.23). 所以速度势方程的解只满足纳维 - 斯托克斯方程而不能满足该方程的定 解条件，因而它不是粘性流动的解.从数学上看，无旋条件使纳维 - 斯托克斯方程退 化为欧拉方程，但退化的方程的解只能满足粘性流动的部分边界条件，即不可穿透条 件.要同时满足纳维-斯托克斯方程和无滑移条件的无旋流动解不存在，这就说明了 粘性流体流动一定是有旋的. (2) 粘性流体运动的起散性 在不可压缩牛顿流体流动的能量方程中有一粘性耗散项 $\phi$ = 2$\nu$ SijS ij/P>O ， 它使 流体质点的情增加，也就是说绝热系统中牛顿流体运动是情增的不可逆耗散系统.

(3) 粘性流体运动的扩散性 粘性流体运动方程 (8.19a) 右端含有二阶偏导数的粘性项，它们具有扩散性质. 粘性流体的扩散性对流动形态的演化有很大影响，比如说，势力场中不可压缩理想流 动中有涡量保持定理(凯尔文定理) :有旋的流体团始终有旋，无旋的流体团则继续 保持元旋.但是，当有粘性存在时，涡量可以通过粘性扩散到无旋区，使得有旋区逐渐 扩大. 以孤立涡管为例，扩散使涡管截面随时间不断扩张.设在势力场作用下，初始时 刻不可压缩流体中有一根具有环量 F 的直线涡，在它周围产生平面圆周运动 z 功 =U=Z

\[L\pi r\]

用纳维-斯托克斯方程研究该流场的涡量随时间变化，由于初始流场的轴对称性，演 化的流场速度将是(用柱坐标) :

\begin{gather*}
V, = 0 , Vo = U (r , t) , v. = 0\\
\omega ,  = 0 ,  \omega . = 0 ,  \omega .-\omega  (r ,  t)
\end{gather*}

对纳维-斯托克斯方程求旋度，可得涡量方程为(由读者自己推导)

\["-h 1-r + z\omega \]

一$\mu$ 吨。-吨。，，，，，，${\rm ‘}$、 $\nu$ 一-

\[h-h\]

它是一个简单扩散问题.方程的解应在 r=O 处有限 (t=O 时除外) ，可得如下形式 的解 z 山 ， t) = \textasciitilde{}乙 exp 卜兰\} Z$\pi$$\nu$ t 飞 4$\nu$ l I 涡量场的扩散示于图 8.4. 70 10 卜.$\leftarrow$一一 飞 60 50 主 40 3 30 20 。 2 3 4 5 r 图 8.4 直涡线的扩散 实线 1= 1. 5 ，细虚线: 1=3.0; 点画线: 1=4.5

从图 8.4 中可以看到涡量沿半径方向是高斯分布，它有 一 当量的半径 $\sigma$ =2 ..ht. 它随时间不断增长，也就是说有旋流动的范围不断扩大.所以，在真实的粘性流体中， 理想的直线祸不能保持线涡，由于粘性，它逐步扩散到全流场.

\section*{8.3  粘性流体运动的相似律}\addcontentsline{toc}{section}{8.3 粘性流体运动的相似律}

纳维-斯托克斯方程 (8.17) 的左端含有非线性的对流加速度项，右端有粘性扩 散项，因此粘性流体的控制方程是非线性的对流扩散方程.除了少数简单流动能获得 解析解外，很难用理论方法获得复杂 三 维流动的精确解;即使用近代高速计算机也难 以获得复杂 三 维流动足够精确的数值解.同此，工程中还常用实验方法研究复杂流 动.例如，常在风洞、水洞或水槽中进行模型实验.如果研究对象是尺度很大的飞机、 船舰等，就不得不用缩小的几何相似模型进行实验.这样就需要知道什么条件下模型 实验的流场能真实地再现实物流场.流动的相似律将讨论该问题. 1.相似定律 (1)几何相似 在初等几何中知道:对应长度成比例、对应角度相等的两平面几何形状称为几 何相似.我们将这种几何相似的概念推广到流动相似. (2) 流动的力学相似 定义 8.3 在时空中几何相似对应点上的物理量成比例的两个流场称为相似 流动. 流动相似的前提是几何相似.在几何相似的流场中，除了实物和模型必须几何相 似外，流场的其他几何边界也必须相似.例如绕流问题的来流攻角必须相等.在几何 相似系统中的相似对应点上有以下等式:王.!.=主.!.=生=旬，立 =$\alpha$ t' 其中 $\alpha$L'$\alpha$，是几

\[X2 Y2 Z 2 --. . 12\]

何相似系统的长度比和时间比，并且全场为相同常数. 流动相似要求以下等式成立:

\[P \alpha  --x-x hr-hy\]

V $\alpha$ 一-

\[x-X JZ\]

E、

\[A-z u-U\]

(8.26 ) 式中的，的分别是速度比和压强比，且在全流场为常数.下标 "1" 和 "2" 表示对应的两 个相似流场. (3) 特征参量与无量纲量 相似流动中某一指定状态的物理量称为特征物理量，在流动问题中应有以下特 征量.

几何特征量 L: 常称特征长度，如绕流物体的直径、机翼的弦长、厚度等，几何相 似系统只需 一 个特征长度; 时间特征量 T: 也称特征时间，如非定常流中的振动频率的倒数，或定常流动中 特征长度与特征速度之比等 z 速度特征量 Uo : 常称特征速度，如定常绕流的来流速度，旋转物体的切向速 度等; 压强特征量 Po: 常称特征压强，如来流压强或滞止压强等; 密度特征量 $\rho$1 :常称特征密度，如来流密度或滞止密度等; 其他物性特征量还有 z 特征粘度 $\mu$。 、特征导热系数 '\textbackslash\{\} 0 以及特征质量力，如重力加 速度 g 等. 物理量与其特征量之比是无量纲纯数.称为无量纲物理量，加上标"一"表示.如 z U, ${\rm Þ}$ cu U; 二 百' $\rho$= 式 . 1 = 亏七等等 它们分别为无量纲速度、无量纲压强和无量纲时间等. 定理 8.1 相似流场中，几何相似点上无量纲量相等. 证明 根据力学相似的定义式 (8.26). 以速度为例.在相似流场中任取两对相似 点的速度比应当相等，即 $\alpha$ 一一 队 -u一-u-u 如果取 U 1 ， U 2 为特征速度，则由此很容易导出 U;) U; 2

\[U BI =U;=U ; =un\]

证毕. 在相似流动中任意物理量 q 和它的特征量 Q 都有类似的关系式$\cdot $旦!..=旦三. QI Q2. 该定理说明 z 在相似流动中只需要有一个无量纲解就可通过无量纲关系得到相 似系统中任意 一 个流动的解，这就是相似理论的重要意义.下面讨论流动相似的充要 条件. (4) 流动相似的充要条件 理论上任意 一 个流动由基本方程和定解条件唯一确定，因此可以由无量纲方程 和边界条件出发导出流动相似的充要条件.以不可压缩流动为例，首先把基本方程和 边界条件写成无量纲形式，有量纲的纳维-斯托克斯方程组是 z i1U; I rr i1U; \_ L' 1 i1 P I J.L ....,2 "2" '+U;7一 =1 一--一 +$\iota$ V'U; dl ${\rm •}$ c1 Xj P dX , P rJU; rJ X;

取特征量 : T( 时间 ) ， L( 长度 ) ,Uo (速度 ) ， g( 作用在流体上的质量力强度)， po 〈压强) ，于是可建立无量纲量 t = 主， $\iota$= 旦 ， Ui = 旦， I， =L ， p= 主U(I' J' g 将它们代入有量纲的纳维-斯托克斯方程，得无量纲的连续方程和运动方程 z (号\}手+(旦 )UEEi=U)J 一(均主主+(码)亨 2 眈I ai 飞 L I\textasciitilde{}) ${\rm e}$J ${\rm I}$) '1'>' J' 飞 $\mu$$\tilde{I}$ 缸，严' (旦)旦旦 =0L lai: \_.-h，，，，"'\_\_l\_.J.IL..$\tilde{Jlt}$l.L... UO \textasciitilde{} ----'..Lr=t nA "， U\textasciitilde{} 式中亨=一寸，是无量纲算符第 式消去常茅数」，第一式同除以」，经整理后，得a-Z $\cdot $一--...."" \textasciitilde{} ". ,.",-\textasciicircum{}" L 如下无量纲方程:

\[au, = 0\]

aXi aU; , YT ${\rm e}$JU; $\tilde{a}$ fj . 1 a2 Ui , 1 Sr 一:::!.+Uj 气!. =-Eu 工$\zeta$+ 一一一:' +:.. f 式中新的无量纲数有 dl . -) a${\rm I}$j -- ${\rm e}$J.ii ${\rm •}$ Re ${\rm e}$J ${\rm I}$jaXj . F, Sr-L. De 一卢。 L \_ UoL 一一 一 …一一一

\[Uo T' ---\mu \nu \]

Fr = 哇 ， Eu = 主可

\[gL\mu )o\]

(8.27) (8.28) 它们分别称为 z 斯特劳哈尔巴trouhal) 数、雷诺( Reynolds) 数、弗劳德 (Froude) 数和 欧拉 (Euler) 数. 流动问题边界条件的无量纲表达式如下. 固壁条件:由固壁无粘条件 Ui=O ， 很容易导出

\[Ui = 0\]

来流条件:由来流条件 (Ui 儿 =Uoco旬i' 也很容易导出(以 U。为特征速度) : (Ui $\lambda$ = co 旬， = const. 界面运动学条件 z 有量纲的界面运动方程为主 +u. V ， =O ，它的无量纲形式是 主一笠 +${\rm Ð}$. 亨 t=OUo T ai ,- ${\rm •}$ \textasciitilde{} 式中 E 是元量纲的界面函数. 界面应力条件:当存在表面张力时，界面两侧的正应力差如下: $\sigma$L 一 $\sigma$.n - $\gamma$( 去+主) 它的无量纲式为

$\sigma$;"-Y/l-L I\} 一 $\sigma$.. - Lpo \textbackslash\{\} 瓦 T R2 \} 式中 y 是表面张力系数，孔，凡是界面的曲率半径. 以上导出的元量纲边界条件可归纳如下. 固壁条件 粘性流体运动的相似律8.3

\[(8.29) Ui = 0\]

来流条件 (8.30) (V,L = co 钮 i = const. L ar 一 \textasciitilde{} .\_ 一一-斗十 U. Vt= 0 Uo T 坷' 界面运动学条件 (8.3 1) (8.32) 界面应力条件 $\sigma$:，， -6-Y/ 土+土\}一一=-- =-- Lp o \textbackslash\{\}R , . R2 I 在几何相似条件下，无量纲边界条件公式 (8.29) 和式 (8. 30) 自然满足，式 (8.3 1) 中无量纲量 dz于是斯特劳哈尔数，式 (8. 32) 中还出现新的无量纲数川=式，称韦 伯 (Weber) 数.于是，元量纲基本方程和边界条件中含有的元量纲参量有 : Sr 、 Eu 、 Re ， Fr 、 We 和 $\alpha$i' 因此该方程组解的一般形式为 U = Ui(.r, ， i. 段 ， Re , Fr .Eu. We. $\alpha$;) $\rho$= ${\rm þ}$<:ii ，[' 段 .Re.Fr.Eu ， We. $\alpha$i) 几何相似系统中 $\alpha$ ，必相等，如果方程和边界条件中所有无量纲参量 : Sr 、 Eu ， Re ， Fr ， We 相等，则对应的流动问题的无量纲方程有相同的解，也就是说流动相似.于是我 们有流动相似的克要条件. 定理 8.2 几何相似的牛顿型流体的流场中，无量纲参数 Sr ， Eu ， Re ， Fr ， We 相 等，则流动相似. 流动相似定理告诉我们，无量纲参数是决定流动相似的重要条件，因此它们常称 为相似准则. 2. 相似准则的物理意义 (1)雷诺数 (Re) 的意义 雷诺数情的是相似流动中惯性力和粘性力量级之比，取相似系统中有代表性的 项分析，就有 nm u叫了州一 一切-h-u-Y-U-a-a M丁一队

\[-u\]

而-h-u-2yu--a 叮

所以雷诺数越大表明流体质点惯性力相对于质点上的粘性作用力也越大;反之 ./J 、雷 诺数流动是粘性主宰的流动.通常用雷诺数对粘性流动分类. (2) 弗劳德数 (Fr) 的意义 弗劳德数 (Fr) 是相似流场中惯性力项和重力项 量 级之比，取代表性项，则有 ，v但应 ${\rm U}$ aD ${\rm e}$J.r 1. a .r 二 U\textasciitilde{} \textasciitilde{} 一 二一 … 一 一… $\rho$'g p g -- gL 卜， 在重力有重要作用的流场中弗劳德数数是关键的相似准则，例如:;.J<波现象主要是 惯性力和重力间的平衡作用. (3) 斯特劳哈尔数 (Sr) 的意义 斯特劳哈尔数 ( Sr) 是相似流场中局部加速度和对流加座度最级之比，取代表 项，则有 aU Uo au 一一一 ---.\_-.\_ $\iota$ 一- .-" u\textasciitilde{}旦 旦 D${\rm e}$J 坐一 U" T ' J '

\[eJ.r L - ax\]

所以 Sr 是流动非定常性的标志，如振荡图性的绕流，圆柱的振蔼周期是重要的特征 参数.对于定常流动，没有特征时间，就元币i Sr 的相似. (4) Eu( 或 Ma\} 和 We 的意义 类似的推导可以证明: Eu 是压强梯度和惯性力间量级之比 ;We 是表面张力和 压强间量级之比. 在可压缩流场 It T 常用 Ma 来代替 Eu. 它们之间有如下的变换式 z Eu 一 $\rho$$\upsilon$ y，$\rho$ u \_ c \textasciitilde{} \_ 1 一二一一一一二 pUi mui yIli yMui 式中Ma。 是特征马赫数，它代表惯性力与压强合力的量级之比(第 7 章已有详细讨论). 了解相似准则的物理意义对于分析流动现象和设计模型实验都是很重要的.下 面说明相似理论的应用. 3. 相似理论的应用 (1)完全相似和部分相似 全部相似准则相等的几何相似流场称完全相似流场.但是实际上在模型实验中 要保证完全相似几乎是不可能的.譬如进行船舰阻力的模型实验时，如果要求流动完 全相似，则 Re 和 Fr 必须相等.设实物和模型的速度和长度特征量分别为 V$\rho$. Lp. V例 .Lm' 如实验在相同介质中进行，物性参数相同.要求 Re 和 Fr 相等，则以下等式 必须同时成立: Re = 岂止主 = 岂止主 I.J\_ P. Fr = 旦i = UL gL" gLm

也就是要求以下等式成立: U$\tilde{U}$;H U.L 命 = U\_L\_. , ,:, p = 一-

\[I'-P - m-", ' l.....\]

L ", 上式只有 一 个解: 1-"=1-，，，和 Up = U"， ， 就是说，模型流场和实物流场完全相同.因此 在相同的介质中进行缩小模型实验时不可能同时满足 Re 和 Fr 相等.如果一定要在 模型实验中保证这两个相似准则同时相等. 一 种可能的方法是采用不同的介质进行 模型实验，例如改变介质的密度或粘度，不过这要付出巨大的代价.在高速气体流动 的模型实验中也有类似情况.用雪气为介质的气动力实验中不可能同时满足 Re 和 Ma 相等. 因此.实际 T 程实验'1'只满足部分相似准则相等.称为部分相似.譬如常常把船 舰的粘性阻力和波阻分别实验.粘性阻力实验中不考虑表面波的影响，只保证雷诺相 似;在波阻实验时，只保证 Fr 相等.而模型 Re 只要求达到 一 定数量，并不要求它和 实物 Re 相等. 还有-血流动 qJ. 当相似准则到达某 \_"值后，流动的性质几乎不随相似准则变 化.例如高 Re 的饨体绕流阻力.当 Re 达到 一 定值\{例如.大于 10 '; \textasciitilde{} 10 7 ) 以后，阻力 系数几乎和 Re 无关.这种现象称为自相似.当流动进入自相似状态，相应的相似准 则就可以不号虑. (2) 相似理论的意义 首先.相似理论告际我们流体运动的科学实验方法.譬如影响低速飞机气动力的 因素有 z 飞行速度 U 、飞机的氏度 L 以及环境的物性参数 :$\rho$ .p . ${\rm •}$ $\mu$. ….假如没有 相似理论，我们必须对每 $\cdot $ 个参数进行实验.而实际 t 只需控制 一 个相似准则 Re. 对 不同 Re 的 T 况进行实验就包括了所有-可能出现的情况.不仅如此.相似理论还告诉 我们:实验结果必须艳理 J;x元量纲系数的形式，这种结果在几何相似的流动 It 1 才有 普遍推广意义. 相似理论推导出的 一 般无量纲关系式还说明:流动的性质并不取决于单独的有 量纲参数.而是依赖于无址纲参数. 一 种给定几何边界的流动.它的流动性质可能随 Re 改变而发牛.很大的变化.例如:简单的民图管 '1 1 粘性流体流动，通常情况下当 Re < 2000-3000 n才流动是有现则的平行自线运动.即通常称之为层流，当 Re 很大时， 如大于 10 丑 ，流动变成很不现则的湍流(第 9 章吓细讨论湍流) "在绕流问题巾也有类 似情况.团性体的定常均生j 绕流'1' .当 Re 1M 小时.流动是规则的(但不同于理想流 动) ;当 Re 达到 100 左右.图柱后出现现则的涡列;继续增加 Re. 涡列开始振荡 ;Re 很大时，阑柱后的流动演变为极不现则的湍流. 总之，流动的无量纲准则是控制流动的参鼓.

\section*{8.4  不可压缩粘性流体的解析解}\addcontentsline{toc}{section}{8.4 不可压缩粘性流体的解析解}

粘性流体的基本方程是对流扩散型的二阶非线性偏微分方程组.到目前为止，还 没有求解该方程组的普遍有效方法.通常有三种途径求解粘性流动问题. (1)简单流动的解析解 在一些简单流动问题中，非线性的惯性项等于零或者可化为非常简单的形式，使 方程组可找到解析解.但是具有精确解的问题为数很少. (2) 近似解 根据问题的特点，略去方程中某些次要项，从而得出近似方程，在某些情形下可 得到近似方程的解.这种途径所得解称为近似解，常用的近似方法之一是参数摄动 \_pUL\_UL 法，粘性流体流动问题主要摄动参数是雷诺数 Re 一一一-一.用雷诺数作为摄动参 数，可以近似求解两类问题: ${\rm 1}$小 Re 问题.这是一类速度极慢、尺度极小或粘度极大的流动. ${\rm 2}$大 Re 问题.这是一类速度较快、尺度较大或粘度很小的流动.直观上粘度很 小的稀流体似乎可以视作理想流体，但是前面已经指出理想流体模型不能求出流动 阻力，但是用大 Re 问题摄动的结果可导出壁面附近的近似粘性流动方程和远离壁 面的理想流动的组合模型，即著名的普朗特边界层理论模型，本书将在第 10 章讲述 这类问题. (3) 数值解 随着计算机的发展，用数值方法直接求解纳维斯拉克斯方程是近年来用于求 解流动问题的一种有效途径.这一方法已发展为专门学科"计算流体力学".本书不介 绍粘性流体流动的数值解法.但是通过解析解和近似解了解粘性流体运动的性质，对 于数值求解方法是有很大帮助的. 下面介绍几个简单层流运动的解析解. 1.圆管内的粘性不可压缩流体的定常层流流动(哈根-泊肃叶流动) 流体介质在圆管中的输送是工程中的常见现象，下面讲述如何用流体力学原理 来提供计算方法和公式. 假定输送管道是水平设置，因此质量力可以忽略不计.如果输送管道距离很长， 则在输送方向上流动是均匀的，若能计算输送距离为 L 的管道上所需的原强差，任 意长度上的压降很容易算出.假定输送流量是恒定的，因此流动是定常的.于是.从流

体力学角度，可求解以下问题. 在无限长水平圆管内的粘性不可压流体的定常层流流动中，已知 z 回管直径 D ， 流量和流体的物性(如:密度和粘度) ，计算 相距长度为 L 的两截面 l 和 2 间的压强差 $\rho$ .-P2( 参见图 8.5). z 下面利用牛顿流体基本方程和边界条件 求解该问题.流动的几何边界是因柱体，因此 选用柱坐标来描述和求解该问题是最合适 的.取固结于圆管的柱坐标系 (r ， 8 ， z) ， 如 图 8.5 所示. 因 8.5 圆管泊肃叶流动示意图 不可压缩牛顿流体的运动方程在柱坐标系中的表达式为 边界条件为 斗 Vr ${\rm •}$ 1 dVn 习" ''\_-r + 一 "\_-: +一二+立 =0ar . r a8 . az . r 生!: +vr 生二十旦旦\textasciitilde{} + v . 主r 't主Jt . -, Jr . r ${\rm e}$J 8 . - , ${\rm e}$J z r 1 J b. I 凹二 2 J$\gamma$18 \textbackslash\{\} =一一 ，: y + $\nu$( 60v r 一-手一 $\tau$ -\_-: )

\[p\theta  r . r -- , r ' r' \theta  8 I\]

旦旦 +V 旦旦+旦旦旦+队旦旦+旦旦2

\[<J t ' -r rJ r ' r eJ8 ' eJ z' r\]

=- \_!\_ \textasciitilde{}\textasciitilde{} + IJ( 60vn + 三生E-21\}pr a8 ' r \textbackslash\{\} --. ${\rm •}$ r${\rm ‘}$ ${\rm e}$J8 r' I 生E 十 V r 旦旦+旦旦旦 +V 豆豆 =-l\_ $\tilde{P}$ + $\nu$601儿J t ' -, a r . r ${\rm e}$J 8 ' -. ${\rm e}$J z P ${\rm e}$1 z ' r - -, (8.33) (8. 34a) (8. 34b) (8. 34c) V r I r= 号 = Vo I r = 号 = V z Ir=\textasciitilde{} = 0 (8.35) 应用流动的基本方程求解问题时，首先应当分析具体问题的流动特征，根据它的 特征简化问题和选用最方便的解法.本问题中，流体在无限长直圆管中由流向压降驱 动.因此流动是单向的平行流动;并且由于回管无限长，单向流动沿流动方向是均匀 的、在周向是轴对称的.根据以上分析，要解的流场可简化为 V. = U( 忡 ${\rm •}$ V r = Vo = 0 ， $\theta$ jJ8 = 0 ， $\theta$ UjJz = 0 (8.36) 分析实际问题，提出简化流动模型，是用流体力学理论解决问题的重要步骤，某种意 义上说，它比求解方法更为关键.简化模型是否正确的依据是能否符合物理实际并满 足基本方程和边界条件，也就是说，如果简化模型和基本方程及边界条件不相矛盾， 这种简化模型就能成立.下面把式 (8.36) 的模型流场依次代人基本方程和边界条件， 可得以下结果.

(1)由于 V r =Vo =0 .JV./Jz=JU( r) / Jz=O. 代入式 (8.33). 连续方程得到满足 z (2) 由于 Vr=Vo=O. 代入式 (8.34a) 得$\theta$$\rho$ /Jr=O; (3) 由于 Vr =vo=o.a pja()=O ， 代人式 (8.34b) ，。方向的运动方程自动满足; 由 (2) 和 (3) 可得压强只是流向坐标的函数，即 $\rho$=$\rho$ ( z). (4) 由于 Vr =Vo =0 ,r7v. / Jr= r7 v. / cJ ()=O. 代入式 (8. 34c) ， z 方向的运动方程简化为 ap = $\mu$ ${\rm ð}$. v = .uI \_!\_且 \{r 亚门= $\iota$二主主=常数 '， z 严严 L r dr 飞 dr J J (8. :37) 式中 Jp / Jz 只是 z 的函数，而 $\mu$ ${\rm ð}$.U 只是 r 的函数，要使等式成立，两项都必须是 常量. (5) 由于以 =Vo=O ， 代人式 (8.3 日，求解方程 (8.37) 的边界条件为

\begin{equation*}U(r) I r c. R = 0\tag{8.38}\end{equation*}

简化的基本方程和边界条件 (8.37) 和 (8. 38) 构成定解问题，只要解出该边值问 题，它就是满足基本方程 (8.33) .方程 (8. 34a)-(8. 34c) 和边界条件 (8. 35) 的解'.由 简化后的运动方程 (8.37) 可看出，非线性的惯性项消失了，只需积分两次，就可得到 它的一般解'.将式 (8.3 7>积分一次得 再积分 一 次得

\[r dU-l d\rho  2 +C,\]

一一一一十. . dr 2$\mu$dz' ，\textasciitilde{}，

\[1 d /J\]

U(r) = \textasciitilde{} \textasciitilde{}互 r 2 十 cllnr + C2

\[4\mu dz' , -,\]

根据该问题的物理特性，在管道中的流动速度应处处有界，所以必须有: CI =0. 否则在圆管中心 r=O 处速度等于负无穷.由管壁边界条件式 (8.38): U(r)lr =R=O. 得 cz=-i 业R 2 ${\rm •}$ 因业=一主二生，速度场的解为

\[4\mu dz" 'F---> dz L\]

1 I PI - P2 \textbackslash\{\} U(r) = $\tilde{I}$ 一一一 1' 2 ) (R 2 \_ r2 ) 4$\mu$ 飞 L I 有了速度分布后，很容易得到体积流量公式: r R 一 $\pi$R 4

\begin{gather*}
(\rho 1 - pz)\\
Q=2 \pi joU(r)dr-8\mu L
\end{gather*}

(8.39) (8.40) 式中 R 是因管半径.式 (8.40) 是因管中层流运动的流量和压降间关系式.下面还可 以进一步讨论该流动的一些主要性质. (1)圆管截面上的速度沿径向是抛物线分布 见式 (8.39 )， (2) 最大速度在 r=O 处， U\_.. = l\_( 且二坐飞R z m.. 4$\mu$ 飞 L I (8. 41)

(3) 平均速度(厕管截面上的平均速度) 利用流量公式，可得平均速度 U\_ = 主$\tau$= 且三生R 2 = $\tilde{U}$\_

\[m\pi R 2 8 \mu L -- 2- m\]

就是说阴管巾平均速度是最大速度之半. (4) 粘度计公式 (8.42) 阑管层流运动的流量公式 (8.40) 由哈根 -7自肃叶( Hagen-Poiseuille) 最先导出，故又 称哈根 - 7自肃叶公式.如已知两截面间压差、管径和粘性系数，就可计算流量.该公式还 可以用于测定流体粘度.设计 一 个很长的阿管流动实验装置，通过测量管径、流量和长 度 L 上的压降，则由流量公式 (8.40) 可得到流体的粘性系数 $\mu$ 的计算公式如下: (5) 沿程阻力系数

\[\mu =\pi  D4( \rho  1- P2)\]

128QL 定义 8.4 管内流动的沿程压降的无量纲 tt 称为阻力系数，用 A 表示

\[\lambda - PI - pz D\]

一 一一

\[rfJ ~/2 L\]

由式 (8.42) 容易导出阅管层流的情程阻力系数为 (8.43 ) 一 $\rho$l 一户2 D \_ 64p \_ 64 A 一一一一-一一一一一一一一 (8. 44) rfJ $\tilde{j}$2 L pUmD Re 式中 Re=pUmDI$\mu$ 是因管内流动的雷诺数. 注意:以上结果对应无限长阿管rJ l 不可压缩牛顿型流体的层流运动，又称"完全 发展的阴管层流流动"若是有限长阿管.则本节公式在管道进出 U 处不适用.关于厕 管进口段的层流运动，粘性流体动力学的专著中有详细讨论.此外.本节给出困管层 流流动的阻力特性，关于回管湍流流动的阻力特性将在第 9 章中给出. 2. 两平行平板间的流动 再研究一个简单的层流流动问题.水平放置的两块无限大平行平极间充满了不 可压缩牛顿流体，平板间的距离为 211 ，如图8.6 所示.已知上板以等速度 U 沿 . 1"铀正方向运动， 下极同定.截面 1 和 l 截面 2 上恒定压强分别为 $\rho$1 和扣，求平极间速度分布及应力分布. 流动的几何边界是平行平面.用直角坐标描 述该流场最合适.坐标系 (x.y.z) 固结于下平板， y=O 位于上下平板的中心平面.此时不可压缩 牛顿流体运动的基本方程为 U f

\[-<:1 I\]

网 8.6 平商库埃特流动示意图

${\rm e}$J u , ${\rm e}$J v , ${\rm e}$Jw -一十一-十一一 =$\sigma$${\rm e}$J:r . ${\rm e}$J y . ${\rm e}$J z (8.45) 旦旦 +u 旦旦 +v\textasciitilde{}旦+吼，但 =-i E主 +II( 主旦+旦旦+主旦) (8.46a) $\theta$ t . -- ${\rm e}$J:r . - ${\rm e}$J y . - ${\rm e}$J z p ${\rm e}$J:r $\nu$\textbackslash\{\} ${\rm e}$J:r2 ' ${\rm e}$J y2 ' (Jz2 I 旦旦 +uE卫 +v 旦旦 +ttJEE=-iE主 +II( 主卫+旦旦+旦旦) (8.46b) 3t3ray az payv 飞 ${\rm e}$J:r 2 ' ${\rm e}$J y 2 ' JZ2 J 旦旦 +u$\tilde{W}$+v$\tilde{W}$+ 悦，但=- \_!\_纽 +II( 生 +21芋+生) (8.46c) ${\rm e}$J t . -- ${\rm e}$J:r . - ${\rm e}$J y . -- ${\rm e}$J z p ${\rm e}$J z$\nu$\textbackslash\{\} J:r2 ' a y2 ' ${\rm e}$J Z2 I 边界条件是 u I y\textasciitilde{} - h = v I 户士 h = w I 户士h = 0 , U Iy $\tilde{h}$ = U (8.47) 根据流动的边界条件，我们可以推测，该流场可简化为

\[u = u(y) , v = w = 0 , eJ /\theta  :r=eJ/eJ z=O\]

即流动是定常的平行平面流动.将设定的流场代人基本方程 (8. 45) 和方程 (8.46a) (8. 46c) ，则连续方程已经满足，并由 y 方向动量方程和 z 方向动量方程可导出 ${\rm e}$J p/ ${\rm e}$J y= ${\rm e}$J p/az=o ， 也就是说，压强只是 z 的函数 .:r 方向的动量方程和边界条件简化如下: 积分式 (8.48) 两次得 ${\rm e}$J 2 U \_ d$\rho$ …一${\rm ‘}$户1 - P2 $\mu$ 巧E 一忑$\tilde{l}$: Un :S l. 二一一$\tau$-

\begin{gather*}
u Iy ~ -h = 0\\
U Iy\tilde{h} = U
\end{gather*}

u = \_!\_兰Ezf+ctv+$\alpha$ $\mu$d:r 2 且 J I ....\textasciitilde{} 一一($\rho$1 - P2 )L2 \_L U 应用边界条件得: Cl=U/ 凯 ， C2 一一一一一${\rm o}$.:- h Z +一，于是所求问题的解为2$\mu$ IJ .- . 2 u= 红毛主 (h 2 - v2 ) +旦(\textasciitilde{} + 1) UL ,.- J" 2 \textbackslash\{\} h . - I 由以上结果可得到该流场有以下性质. (8.48) (8.49a) (8.49b) (8.50) (1)它由两部分线性叠加组成，一部分是压降驱动的流动，速度是抛物线分布 z 生二主主 (h 2 -y2); 另一部分由上平板拖动，速度呈线性分布$\cdot $旦\{主 +1 ). 2$\mu$L ....... J '7/ <1 ........"'<1..............1...- I -.，.....，r.g-:T<l 7\textasciitilde{}/..""..\_， ..::::.:::::;.-........I........ \_，.I - I t". 2 \textbackslash\{\}h ''''J (2) 剪应力分布 应用牛顿剪应力公式可得

\[du - Pl - P2--\]

xy - p. $\tau$一一一一一:- P2v +丘U${\rm •}$ dy L - 2h (8.51) 式 (8.5 1)表明一部分剪应力由压降引起，呈线性分布;另一部分由上板拖动所引起， 剪应力为常数.

(3) 流量公式 z 两平板间单位宽度的体积流量为 fh -' \_ 2 I 如 - P2 \textbackslash\{\},\_ 3 Q = I udy = 芒 l 一一「一 Jh 3

\[+Uh\]

J - h "$\mu$ 飞 L I (4) 截面平均速度 Um-Q-l/$\rho$1 一 $\rho$2 \textbackslash\{\},\_ 2 I U 一一一 $\tilde{I}$ 一一一一 Ih" + :: 3$\mu$ 飞 L I (5) 平面泊肃叶流动 两固定的平行平板间由压强差驱动的流动称平面泊肃叶流动，以上结果中令 U 等于零，得平面泊肃叶流动的速度分布及流动特性如下 z \_ PI - P2/L2 u 一-一一一 (h 2

\[- V2 ) 2\mu L o-- J\]

即，平面泊肃叶流动的速度剖面是抛物线，最大速度在 y=O 处 z u 一 $\rho$ 1\_ -\_P2h 2 一-

\[max 2\mu L\]

流量 Q= 主(乓主 )h 3 平均速度

\begin{gather*}
U~ = PI- -- P2h2 = : u\\
m 3\mu  L'- 3 -max
\end{gather*}

最大剪应力在平板上 (y= 的: 'n hy 一 二 L

\[hr-\]

一-r 沿程阻力系数 $\lambda$ 一 ($\rho$1 - P2)h \_ 6 一 一一

\begin{equation*}pU\tilde{LI}2 Re\tag{8.52}\end{equation*}

式中 Re=pUmhl$\mu$. 3. 库埃特 (Couette) 流动 元限长同心圆柱和圆筒间充满不可压缩牛顿流体，内 柱以等角速度绕轴旋转，这时在环形空间内的流动称为库 埃特流动.下面求圆柱环内的流体速度分布和作用在柱面 上的剪应力. 根据该流场的几何特征，用柱坐标描述该问题最为合 适，使柱坐标的轴线和同心圆柱的轴线重合，如图 8.7 所示. 已知: d l ， d2 分别为内圆柱的外径和外圆筒的内径，图 8.7 库埃特流动示意图

内同柱以 $\omega$ 等角速度转动.由边界条件的轴对称性和驱动条件的恒定性，我们可以 推测流场是定常轴对称的，即 v. = v. = 0 , Vo = v( r\} ${\rm •}$ a I J(\} = 0 将以上速度和压强分布代入式 (8.33)- 式 (8. 34c). 连续方程自动满足 z 由轴向动量方 程可导出 $\theta$ pjaz=O ， 因此压强只是 r 的函数 z 周向动量方程和径向动量方程简化如下 z 该流动的边界条件是 d2 v , 1 dv v \textasciitilde{} --\textasciitilde{}一一一一

\[dr2 ' r dr r 2 -\]

业=但l dr r 七 I . $\tilde{Rl}$ = v(R 1 ) = 护1

\[Vo 1. \tilde{R}2 = v(R 2  = 0\]

设 v=r" 代入方程 (8.53) ，可得 : n(n 一l) +n 一 1=0 ，解得 n= 士 1 .即 v(r) = $\iota$ 1r+ 生 r (8.53) (8.54) (8. 55) 利用边界条件 (8. 55) 求出积分常数 : C1 = -R\textasciitilde{} $\omega$ I (R\textasciitilde{} - Ri> ${\rm •}$ ("2 = - m R\textasciitilde{}$\omega$ I(m - R t>.最后得

\[Ri\omega  (RUr- r)\]

。 = v(r) = 一一'R\textasciitilde{} -Ri 由速度场可求出流场中的剪应力分布 z ,c1v v 飞 n R$\tilde{R}$\textasciitilde{} $\omega$l <.8 -叫百二- -;: ) =- t; $\mu$ 豆7士王;F (8.56) (8.57 ) 流体作用在内柱面上的切应力和旋转方向相反，所以阻力矩也和内圆柱旋转方向相 反，其大小为 f2. , , n' 'n 4'/[J\_l R$\tilde{R}$$\tilde{M}$.=| (r a) u R?dO=J坦一 \_.$\omega$Jo

\begin{gather*}
'.,-u'r-nl--I-'"\\
R~-Ri
\end{gather*}

(8.58) 式 (8.58) 也可用作测量流体粘度的公式，只要测定内回柱上流体作用力矩和转 速以及内外图柱的半径，就可由该式计算流体动力粘度系数. 将速度分布公式 (8. 56) 代入式 (8.54) .然后积分求出，可求出圆筒内的压强分 布.压强的定解条件是必须给定流场一点的压强，例如内圆柱或外圆筒面上的压强. 4. 圆管内幕律流体的定常流动 最后研究在无限长圆管内罪律型非牛顿液体的流动，流动的边界条件和牛顿流 体的圆管流动的情况相同.流体是平行直线运动，即速度场为

\begin{equation*}v. = Vo = 0 , v. = v( r)\tag{8.59}\end{equation*}

${\rm i}$ 1 ，\textasciitilde{}\textasciitilde{}， ?l2 \_L\textasciitilde{}="-\_\_\_'- \textasciitilde{} 1 d 根据幕律流体的流变性哨 =kl2 tr(2S) 吁 ，对于平行流动， Srz =$\tau$ 王，其他分 量都等于零，因此 $\mu$= kl 吉 l' 1 流场中的偏应力张量 z

\[r= \cdots  = '.0 = \eta =0\]

'" = k 1 叮 l 生drl dr 将速度场代人粘性流体运动方程得到 一纽$\downarrow$主立=() 旦主=纽=() rJz . ${\rm e}$J r riJ(J ${\rm e}$1 r 将军律流体的本构方程代入后得 业=立 (bl 生. I dv\textbackslash\{\}

\[e1 z e1 r   \tilde{I} dr I dr J\]

上式成立的必要条件是等式两边为常数，令 dp/dz=C 1 ，则 |du 旷 1 dv h 一|一 = C1r+ C2 drl dr (8.60) (8.6 1) 由流动的对称性可得，在 r=O 处，应有 dv/dr=O ，因而积分常数 (i =0 ， 于是上式可写作 z | 业判叮1 0 1 坐$\iota$C1 dr 1 dr k 回管内由压强驱动的流动中， dp/dz<O ，即应有 C1<0 ，因而 dv/dr<O. 于是上式可 写成 z (古)=一( I 引 r ) 1'. 对 r 积分，并利用管壁无滑移条件得 r=O 处的最大速度为 (r) = $\tilde{I}$ 旦|飞 R(n i ll 闯一 r(n t-l l ; o)

\[n + 11 k 1\]

v\_.. = $\tilde{I}$ 旦 |IMD 〈 ml lvn 一一，-， .,. n + 11 k 1 一 于是事律流体在回管中流动的速度分布为

\begin{equation*}v = vm.. [l- (r/R)' +I/'J\tag{8.62}\end{equation*}

当 n=l 时 ， v=vm..[l 一 (r/R)2J ， 它就是牛顿流体在回管中的流动.图 8. 8 给出不同 事指数时非牛顿液体的速度剖面.当 n<l 时，即剪切变稀液体，其速度剖面比较饱 满，反之 ， n>l 的剪切变稠液体的速度剖面陡峭.

2 0.6 三$\zeta$ 0.4 0.2

\[-1 -0.5 o\]

rlR 0.5 图8.8 幕律流体在回管内的定常流动 第 8 章粘性流体力学基础 -: n=l. 牛顿流体; 0--0: n=2 ， 事律流体; 0--0: n=O.4 ，幕律流体

\section*{8.5  小雷诺数粘性流体运动的近似解}\addcontentsline{toc}{section}{8.5 小雷诺数粘性流体运动的近似解}

对于复杂的粘性流体绕流问题，一般难以求得解析解，除了数值解外，常用摄动 展开法求近似解. 1.正则参数摄动 通常在流动问题中含有几何的或物理的无量纲参数，例如:l./<披问题中的波陡 (波高除以波长)、翼型绕流问题中的相对厚度、在粘性绕流问题中 Re 等是重要的几 何或物理无量纲参数. 如果控制流动的元量纲量是小量，则受数学分析中泰勒展开式的启发，含小参数 (E<< l)方程的解常用小参数作摄动展开，例如

\[Ui (X" E) = uiO' (Xi ,O) + EUi 1l (Xi , O) + E 2 UFl (Xi , O) +\cdots \]

式中按小参量 $\epsilon$ 展开的解 ur 肘 .Ui ll .U: 2l 等可以作为 U i 的各阶近似解，如果这一解 在全流场一致收敛，称该摄动为正则摄动，可得到正则摄动的逐级近似解.倘若该展 开式在流场内非一致收敛时，需要采用其他摄动展开法，如奇异摄动法等. 2. 小雷诺数流动的零阶正则摄动 一一 斯托克斯方程 现在研究 Re 很小的粘性流体运动时，可以用 Re 作小参数展开，令 $\epsilon$= Re<< 1.将 Ui=Ui ( 句 ， Re\} ${\rm •}$ $\rho$=$\rho$ 〈鸟 ， Re\} 对小参数 Re 展开 z

U(Xi' $\epsilon$) = U( 的 (X ，) + EU(l) (X ,) + O(E2 ) \_ p(O) (Xi) $\rho$ (Xi' $\epsilon$) 一一-F」一+ p(l) (Xi) + O($\epsilon$) 以上展开式中，压强的首项比速度展开式的首项高一个量级，这是因为在低雷诺 数流动中，粘性项比压强梯度项高一个量级(参见纳维-斯托克斯方程和下面的分 析) .令 Re=E. 定常流动的纳维-斯托克斯方程可写作:

\[v. U = . \]

因J 1 言- -一 $\tilde{v}$$\rho$ +\textasciitilde{} ${\rm ð}$.U

\section*{1.11  P E}\addcontentsline{toc}{section}{1.11 P E}

将展开式代人无量纲化的定常纳维-斯托克斯方程(不计质量力) .并按小量 E 的 幕次整理，可得 2盯 (0) a盯 11\} ， T牛一 +E .':-+ O(E") = 0 (IX , cJX , l 华主 1 乒fL=u:$\omega$ 掣主+学二十三些L

\[E rJ X ,  \epsilon  dXjdX,  dX, dX , (IX"IXj\]

+E( 学士 +U$\tilde{O}$) 辈二十 Ujl) 型主+…忡。(E 2 ) 飞 (IX ， (IX , (IXj I 将同量级项归并，得各阶近似的线性方程.在零阶近似中非线性惯性项已经略去.事 实上 .Re<<l 意味着粘性力远大于惯性力，因此，在零阶近似中惯性项是高阶小量.在 一阶近似中(运动方程右端 h 则可包含部分惯性效应.小雷诺数流动的零阶近似称为 斯托克斯 (Stokes) 近似，当不计质量力时，零阶近似方程有如下形式的斯托克斯方程 (恢复到有量纲形式) :

\begin{gather*}
v. U = 0\\
V\rho =\mu  \partial .U
\end{gather*}

小 Re 的流动意味着流动速度很小或粘性系数很大，因此斯托克斯近似的流动也称 为极慢运动. 3. 小雷诺数圆球绕流 第 5 章研究过均匀来流绕嗣球的理想流动，现在计算小雷诺数均匀来流绕圆球 的定常粘性流动问题.设无穷远来流速度与 x 轴 平行为 u. i. 圆球直径为 D( 半径为 a). 绕流的 坠 Re=V ." DI $\nu$<< 1，求流场速度分布、压强分布和圆 球所受阻力. (1)采用球坐标系描述该问题，将原点放在球 一+ 心，并使 ()=o 的轴线与来流方向重合. 流动的边界条件是轴对称的，因此可以推断 - ,\_

\[(R. \theta )\]

x 流场也是轴对称的(参见图 8.9) .因而有 图8.9 绕圆球的斯托克斯流动示意图

U o = ve(R ,0) (UR=w 川) , U" = 0 ， $\rho$=$\rho$ (R ， 的 将它们代入球坐标中的斯托克斯方程，得 连续方程 运动方程 ${\rm e}$J v" , 1 av. , 2v" , V峨 cotO \textasciitilde{} 一」二十一一-十--二十」一一一 =0

\begin{gather*}
aR ' R ao ' R 'R -\\
\theta  b ,  a
\end{gather*}

2 引 " , 1 J2 V" , 2 ${\rm e}$J凹 n 立 =u $\leftarrow$一子+一$\tau$ 一一子+一一-"aR ，...飞<7 R Z ' R ${\rm ‘}$ aoz ' R aR 十 cotOJUR2Juo2UH2cotol 一一二一一一\_\_\_"7'\_'

\[R 2 JO R 2 <70 R 2 [i2\]

V O\} 1 纽-..(主旦$\downarrow$土主旦 R ao ,... \textbackslash\{\} a R 2 ' R 2 a6严 第 8 章 2 Jvo , cotO Jvo , 2 dVR Vo \textbackslash\{\} +一-:\_\textasciitilde{}+一一.， .， u 十 IJ .\textasciitilde{} "J( 一一. -." ., - I

\[R\theta R ' R 2 <10 ' R2 <7 0 R2~;j nZOI\]

边界条件 ${\rm 1}$球面无滑移条件

\[R = a , VR = Vo = 0\]

${\rm 2}$无穷远来流条件

\[R \rightarrow \infty  ,  VR = U '. cosO, Vo =- U .. sinO ,  \rho =\rho  , \]

粘性流体力学基础 (8.63) (8. 64) (8. 65) (8. 66) (8.67) (8. 68) (2) 求解 z 用分离变量法解该边值问题.由边界条件可推断解的形式应为 VR = f(R)F( 的 ， Vo = g(R)G (O) ， $\rho$ = ph (R) H(O) + $\rho$ 利用边界条件，确定出 F( 们， G (0) 函数的具体表达式，为了满足式 (8. 6 7)和 式 (8.68) ,F(O) =cosO ,G (O) = sinO ， 即奋

\[VR = f(R)cosO , V(I =- g(R)sinO\]

和边界条件

\[f( \infty ) = g( \infty  ) = U. , f(a) = g(a) = 0\]

将分离变量式 (8.69) 代入基本方程 (8.64 )、方程 (8.65) 和方程 (8.66) 可得 叫[1' + 去 (f 一 $\rho$J= 0 H(O)h' (R) = cosOI f' + 主j'-40-g 门I J 'R J R2 ${\rm •}$ J l">' I h . \_r " 2 , 2 , p , l H'(O) 工= sinOI- g 一 ng 一寸 (I - g) I L '"' R'"' RZ '.' '"'. J 由以上三式，可推得 H(O)=cosO 和以下的常微分方程组 z 1'+ 去 (f-g)=O (8.69a)

2 \textasciitilde{} ， 4 h'(R) = f' + ';.， 1 一 -$\tau$ (I- g) R J R\textasciitilde{} (8.69b)

\[r" ,  2 , ,  2 , ~ , l h=RI g" +\tilde{g}'+~2(I-g) I (8.69c) 1'" , RO 'R2 'J ,." I\]

应用无穷远压强边界条件，可得 I!( $\infty$ )=0. 因此函数 I(R) 、 g(R) 和 h(R) 的边界 条件是

\begin{gather*}
I( \infty  ) = U . . g( \infty  ) = U\\
I(a) = O. g(a) = O. h( \infty  ) = 0
\end{gather*}

求解以上方程组的具体方法如下，先消去 g ， h ， 得到111 的常微分方程为

\[R'/" + 8R\]

3 广 + 8R2 f' - 8 町'=0 该常微分方程的一般解是 I(R) = AR :1 + BR I 十 C+DR 2 由式 (8.69a) 和式 (8.69c) ，可得 A\textasciitilde{} .， B\textasciitilde{} 1 J! (R) = 一': R :i + ':: R I + C + 2DR 2 ..,.--. 2--' 2

\[h(R) = BR 2 + 10DR\]

将边界条件代入以上方程，可确定常数 A ， B.C ， D 如下 z a :i ${\rm ••}$ 3U , a A= 一一一， B =一一-一 ${\rm •}$ C=U. , D=O 2 ,- 2 最后，得 I(R) 、 g(R) 、 h(R) 的解等于 f = ll. 旷 a :l 3 U . a =一-一一一一一 +U2 R i 2 R U. 旷 3 U 、 a ，.. . 3U . a g= 一一一一一一一一一 +U.. ， h = 一一寸-4 Ri 4 R ,\_ ..," 2 R 2 代国分离变 :ht式，得该流场的解为

\[1. 3 a , 1 a"  \]

H = U ., cos8 ( 1 -一一+ :一丁 i\textbackslash\{\}' 2R'2R i l (8.70a) 旷

\begin{gather*}
-r\\
l-4\cdot  u-puo
\end{gather*}

每U

\[3-4\cdot \]

仆 UF AU$\mu$ $\cdot $川

\[3-2\]

U 一 -p 一一一-n叫

\[\rho a\]

(8.70b) (8.70c) 为了求得回球的流动阻力，先求回球表面的应力分量.根据牛顿流体的本构关系 式.可求得球面上的应力分布为 (P HH )H\textasciitilde{} . ， 二(- P + LHH )H =" I ., n rJVH \textbackslash\{\} , 3 /.l U . =(一 $\rho$+2$\mu$ 一:-: ) =一 $\rho$ +一巳二二-cos

\[  r' -,.,. rJ R I H " r" 2 a\]

\[(PRIJ)R\tilde{u} =(TRIJ)R =u\]

\{ 1 iJVR I  Vo Vo \textbackslash\{\} =，一一+一一一'

\begin{gather*}
  R eJ( 'aR R J R \tilde{u}\\
(TO", )R~. = 0
\end{gather*}

=- \textasciitilde{} \textasciitilde{}气in8

\[L\alpha \]

由于流动是铀对称的，流体作用力的合力方向应与来流方向相同，即只有阻力，而无 升力.阻力的大小是

\[D= F-\]

= f: [ (P RR ) R $\tilde{u}$ C叫一 (PRIJ )R川川=勺= u SI吕叮i叫2$\pi$d仇s创叫i =6$\pi$$\mu$ U， .川， a ($\omega$8.7门1 ) 上式称为斯托克斯阻力公式，由此可见，因球所受阻力的大小与来流速度、圆球半径 和粘性系数成正比.常用无量纲阻力系数表示 三 维绕流物体的阻力，它等于阻力除以 plf., /2 与迎风面积的乘积，并用 C$\upsilon$ 表示. 绕圆球小雷诺数流动的阻力系数为 C2D24 一- u-F5歹 - Re (8. 72) 式中 Re 的定义是 Re=U 的 D/$\nu$.D 是圆球直径.根据近似展开的条件，上式适用于 Re<< l 的圆球阻力，例如:非常细小的液滴在空气中的等速运动，或细小的粉末在帖 性流体中的运动等.该流场的流动图案如图 8.10(a) 所示(困球向左运动) . (a) 公正L 歹有军 (b) 图8. 10 绕网球小雷诺数流动的流场 (a) 斯托克斯近1'I:J. (b) 奥辛近似 (3) 结果分析及奥辛 (Oseen) 修正 上面求解了纳维-斯托克斯方程的零阶近似，该近似的基础是流动质点的惯性 项远远小于粘性表面力的合力，下面需要验证上面所得的解是否在整个流场内都满 足近似的条件〈一致收敛) .为此，考察所略掉的惯性项量级的大小.例如:考察 r 轴 上 R>>a 的点，由解出的速度场可得到这些点上流体质点的加速度为

\[D VR I eJ VR I 3 U~.. a  1 a2    1 3 a I a3  \]

Dt l 凹 = VR ${\rm e}$JR 10=0 = "2 V \textasciitilde{} 1 - R2 \} \textasciitilde{} 1 一言 R T 2R3 \} 单位质量粘性力的大小与 if主同量级，即pa.n

\section*{8.5  小雷诺数粘性流体运动的近似解}\addcontentsline{toc}{section}{8.5 小雷诺数粘性流体运动的近似解}

1 a 主 I =坠旦J

\[p\theta R I O\tilde{O} pR:1\]

因此惯性力与粘性力量级之比是 惯性力一 U 、 R \{1 一豆\textbackslash\{\} \{1 \_主旦..L工 u 福E万一丁;-\textbackslash\{\}$\gamma$A 豆豆/川\textbackslash\{\}' 2R I 2RlJ 由此式可看出，随着 R 的增大.此比值近似为 U 民 R/2 川当 R 趋于无穷大时，比值趋于 无穷大，可见在远离阔球的地方，总会出现惯性力超过粘性力的情况，在这些区域内， 所得解与假设条件"惯性力量级远远小于粘性力量级"相矛盾，因此斯托克斯近似所 得流场的解只适用于物体附近的区域. 为了提高近似的准确度，奥辛于 1910 年对小球在粘性流体中运动问题提出了 "在运动方程中保留主要惯性项，略掉次要惯性项"的修正意见，得到了奥辛近似解. 他将速度写成

\[U = U .i+ v'\]

其中 u' 是附加小量，保留一阶小量，质点加速度可近似为 普= [(U ... i+ 山. V]U 句 U 空d:r 于是纳维-斯托克斯方程可近似为  U 1 \textasciitilde{}. , /A $\tau$$\gamma$= 一 ..:\_v$\rho$ +.!::.. ${\rm ð}$.U o :r p p

\[v. U = 0\]

式 (8.73) 和式 (8. 74) 是奥辛修正方程，相应的边界条件为 R $\rightarrow$ $\infty$ :U =U 、 = U. i (8.73) (8. 74) (8.75)

\begin{equation*}R=a:U=O\tag{8.76}\end{equation*}

与斯托克斯方程相比较，运动方程中增加了部分惯性项.在小球附近区域中，由于惯 性项比粘性项小得多，奥辛近似的解和斯托克斯近似解相差很小 P 在足够远的地方， U ${\rm ‘}$ aU/ r1:r 是惯性力的主要部分，保留这一项可以提高解的精度.在相同边界条件下 奥辛得到了比斯托克斯精度高的一阶近似解(具体解法从略) .奥辛近似的流场示于 图 8.10(b). 奥辛近似的回球绕流阻力: D= 瓦 =6 叩U.a(l+ 哈) (8.77) 奥辛近似的阻力系数 z CIl =坐 (1 +立Re)" Re 飞 16- -- / (8.78) 式中 Re=U. D/$\nu$.D 是因球直径. 图 8. 11 是小雷诺数网球绕流阻力系数的斯托克斯和奥辛近似的计算结果，在 Re<<l 时，斯托克斯近似是很好的 ， Re注 1 时斯托克斯近似的计算结果明显偏小.

\[?qUM-\]

理论

\[r 24 1 rl L/)= R~\]

24/..3"\textbackslash\{\} 11 C/)= 一IH ;=7 R...1R... 飞 16'''' HM 近赔似克近托辛验斯阳或实IHA

\[-2 -1 0\]

IgRe 图 8. 11 小雷诺数的斯托克斯近似初奥宇近似 4. 润滑问题 小雷诺数近似可以应用于轴承润滑问题.大型高速旋转机械常用滑动轴承以减 少摩擦和震动.滑动轴承采用粘度较高的润滑油，轴和轴套之间的间隙 $\delta$ 很小，虽然 旋转轴的周向速度比较大，流动的特征雷诺数 U8/v 仍然很小.因此润滑问题可以采 用斯托克斯方程作近似分析和计算.为了分析简明起见.我们把轴承的回周方向展成 平面，分析如图 8.12 所示的间隙很小的平面模形中的二维定常粘性流体运动.

\begin{gather*}
nunu zlA-\\
E-n
\end{gather*}

吨 。龟 l 龟 0.8 1.0 X 一- -1 图8.12 平面润滑问题的 ，;i\textasciitilde{} 意图 上阁 z 润滑面上的压强分布; F 因 z 平面轴承内的速度分布 在直角坐标中，定常二维斯托克斯方程可写作 dU ， $\lambda$ 啊 二十:::....:: = 0 ${\rm e}$J:r ' ${\rm e}$J y

\[1 d IJ , I eJ 2 U ' eJ 2 U  \]

一一二E 十川一$\tau$ 十一一亨 )=$\sigma$ P ${\rm e}$J:r ' ${\rm •}$ \textbackslash\{\} iJ:r' ${\rm e}$J y" I (8. 79) (8.80a)

1 J主 +$\nu$( :2 与+主导 )=0 (8.80b) P Jy 飞 \{\}x" ' (\}y" J 润滑问题中间隙远远小于轴承的周长，即 h<<L. 因此 h/L<< l.在 y 方向用间隙 h 无量纲化， .1'方向用 L 元量纲化，即 x=Lx , y=h y ， 于是有 a a a r1 石 -PZ' 巧 E3 元量纲导数应当属于同一量级，由上式很容易得到 主》立和c1 y -- ax 将以上结果代入式 (8.79) ，可以得到 <1 2 ${\rm ••}$ <1 2

\begin{gather*}
-- ..-r1 y2 // <1x 2\\
u>> v
\end{gather*}

代入式 (8.80a) ，可得 1 ${\rm e}$J 主-"旦旦

\[P r1.1' ~\theta  y2\]

由于 v<<u ， 将式 (8.80a) 和式 (8.80b)1i故 it 级比较，应有

\[ap // c1 \rho \]

一…- rJy "c1 x 即垂直于轴承面的压强梯度可以忽略不计.于是，求解润滑问题的方程简化为

\begin{gather*}
u\tau \\
nuz-
\end{gather*}

飞 d =J$\nu$J m-hd= +业 h

\[h-hl-p\]

(8.8 1) (8.82) 边界条件为

\begin{equation*}y = O:u = U; y = h(x):u = 0\tag{8.83}\end{equation*}

式 (8.82) 有解的必要条件是方程两边等于同一常数.具体求解过程如下 z 先将 方程 (8.82) 沿 y 方向权分，得 1 db ? u= 一旦v 2

\[+ l', v + l'. 2\mu d.rJ 1'" I J I ....r.\]

代入边界条件式 (8.83) ，并简化后，得 1 db. ? u= 一兰左 (y' -hy) + 一 (h - y)

\[2\mu dx\]

设通过轴承的流量为 Q. Q = r u(y)dy = 一生工业+哑J 0 -, J" -J 121-.1 dx' 2 由上式可得沿铀展的压强梯度: 业=但旦一旦PQ dx h 2 h 3

轴承面上的压强分布: f'/6 /J. U 12l$\rho$ 、 $\rho$= II $\rightarrow$亏 一 -4二 )d .r + ( J u 飞 h' h" I 已知模形的倾斜解 $\alpha$. 则 h=hl- .r tan$\alpha$ ，代入上式并积分得 $\rho$ 一年U 6,\_t:J :;-+ -寸-c

\[tan\alpha  (h l - .rtana) tan\alpha  (h l - .rtana)Z\]

最后利用压强的边界条件确定常数 c 和 Q. 轴承内的压强是由模形运动对流体挤压 产生，进入轴承和流出轴承的压强应当相等，如果上式中的压强表示相对于间隙外部 的压强，则压强的边界条件可写作=

\[.1" = 0 ,  \rho =0; .T =L ,  \rho =0\]

将压强边界条件代人压强积分式后，得积分常数 最后，轴承间隙中压强分布为 Q Uh llzz

\begin{gather*}
11 1 +h2\\
c =- 6I.lU\\
tana(h l + h 2 )\\
\rho 6\mu U(h - h 2 ).r\\
h 2( !z I+h 2)
\end{gather*}

压强分布沿 z 方向职分.可得流体作用在轴上的合力丑，即轴承的承载能力 z (8.84)

\[("" 6 f1 UL~ /, hl .,1I1 -h2  \]

)' = I$\rho$d.T =一-乙一一-:-:;-Iln 一一 2 一一一'; z 1 (8.85) J or --' (h l -11 2 ) 2 飞 11 2 - 11 1 十 11 2 I 式中 h l 是轴承进口的间隙高度，润滑油进入轴承后，间隙高度减小，即 11 1 -h 2 >0. 由 轴承承载公式可见 z 承载力与润滑油的粘度和轴承的滑动速度成正比 p 与轴承的受 力氏度平方成正比.此外式 (8. 85) 还表明，润滑油的厚度(常称油膜厚度)越小轴承的 承载能力越大.旋转轴上受到润滑油的摩擦力，摩擦力将消耗一部分能量.由润滑油 速度分布，可以计算作用在轴上的剪应力(请读者自己完成以下两个公式的推导) : 4$\mu$ U 6$\mu$V h,h ( r.", ) 户。= 一一--一一一+ 一'l'L2 一 (8. 86) ","" -" h]- .T tana (h]- .T tan$\alpha$ ) 2 hl+h 2 f" f \textbackslash\{\}.J \_\_ \_ 2$\mu$UL /'>1\_ h l ., h] -h2 \textbackslash\{\} F., = 一 I (rI v ) o d .T一一一一- (21n 一一 3 一一一':' 1 (8.8 7) J 0 --IY 'V -- h 1 - h 2 飞 h 2 - h 1 + h 2 J 上式表明:润滑油膜厚度越小，摩擦力越大，也就是说能使承载力增加的因素，同样 也使铀承的摩擦力增加，因此.在轴承设计时，要综合考虑，做优化选择. 以上讨论的粘性流体运动都属于低雷诺数层流运动的解析解，或很小雷诺数流动 的近似解.当流动特征雷诺数较高时，流动将由层流向湍流过渡，第 9 章将讲述湍流的 发生和充分发展湍流的统计规律.大雷诺数粘性绕流的近似解也不同于小雷诺数绕流 的斯托克斯近似，大雷诺数粘性绕流在物面形成边界层，粘性主要在边界层中控制流体 运动，而边界层外，流动可以视作无粘的.第 10 章将讲述边界层理论和近似计算方法.

\section*{习题}\addcontentsline{toc}{section}{习题}

\noindent\textbf{8.1.} 证明 z 静止物面上具有下述流动的粘性流体物面条件: D ${\rm •}$ n=O. 式中，

D=V XU 是旋度向量 .n 是物面的法向量.

\noindent\textbf{8.2.} 证明:不可压缩牛顿流体定常流动中，当质量力有势时，沿流线坐标 s 有 z

(UJ)(U2 A 一一一 一+主 +II)= $\beta$ 气。是涡量 .II 为质量力势)$\nu$ rJ s) \textbackslash\{\} 2 ' P

\noindent\textbf{8.3.} 半径分别为"和 b 的同心圆柱和阿筒，它们分别以 $\omega$I 和 $\omega$2 作等速转动，因

柱环形空间中充满不可压缩牛顿流体，现给出下列三种情况的速度 z

\begin{gather*}
\rho  '2 b2 wla 2\\
(1) \omega I=zf- ,  uo=-7 ,  ur=03\\
(2)\omega 1=\omega 2. Vd= \omega  Ir.v,  =O;
\end{gather*}

2w句 b 2 (3)$\omega$1 二亏一 ， VO= $\omega$1 r ,v , =0 a 请分别讨论 z (1)流场是否粘性流动的解? (2) 流场是否无粘流动的解? (3) 流场是否有旋的?

\noindent\textbf{8.4.} 油在水平厕管内作定常层流运动，已知: d=75mm.Q=7L/s.p=

800kg/m :! ，壁面上 r毡， =48N/ 时，求油的粘性系数. 8.S 直径为1. 5mm 、质量为 13.7mg 的钢球，在一个盛有袖的直桶中垂直等速 下降，在 56s 内下落 500mm. 油的比重为 O. 屿，因桶的直径和长度足够大，可以忽略 厕桶的端部和壁面效应，求油的粘性系数.

\noindent\textbf{8.6.} 在两块无限大的平行平板间充满两种互不相混的不可压缩牛顿流体，上层

流体的厚度为仇，粘度为 $\mu$2 ;下层流体的厚度为丸，粘度为 $\mu$1. 上板面向右作匀速直 线运动，速度为 U ， 下板面固定，全流场压强是常数，并不计质量力，求流场的速度及 剪应力分布.

\noindent\textbf{8.7.} 证明不可压缩牛顿流体的二维流动，在不计外力的情况下，流函数?满足

下列方程 z 主(，${\rm ð}$.1${\rm þ}$') +丝主(，${\rm ð}$.1${\rm þ}$') - \textasciitilde{}主主(，${\rm ð}$.1${\rm þ}$') = 1I,${\rm ð}$.2 1${\rm þ}$' cJ[ dy d :r d :r dy

\noindent\textbf{8.8.} 设不可压缩牛顿流体在椭圆形截面的无限长直管中由压强驱动，椭圆的长

短袖分别为 a 和 b ， 轴向压力梯度 d$\rho$/d:r是常数，且不随时间变化，不计外力，求该流 动的速度分布和平均速度.

\noindent\textbf{8.9.} 设不可压缩牛顿流体在无限民同铀图柱面间由轴向压力梯度驱动.已知两

回柱的半径分别为"和 \{J. 轴向压力梯度是常数且不随时间变化，不计外力.求该流 动的速度分布和管壁上的粘性剪应力.

\noindent\textbf{8.10.} 半径为"的小球在粘度很大的不可陌.缩牛顿流体中下落，证明其最终速度为

\[U, = 2a2 (p, - p)g/9\mu \]

p 和 $\mu$ 分别是流体的密度和粘度， $\rho$，是小球的质量密度-

\noindent\textbf{8.11.} 不可压缩牛顿流体在夹角为 2$\alpha$ 的两无限大平板间作极慢流动(图 8. 13) ,

若单位宽度的体和、流量为 Q ，外力不计，求该流动的速度分布和 ffi 强分布. ${\rm ‘}$可二Z$\alpha$ 图8. B 题 8. 11 ，示忠阁

\noindent\textbf{8.12.} 在无界的不可压缩牛顿流体中有 一 半径 a 的小球，当小球极慢地绕定轴以

等角速度 $\omega$ 转动时，求作用在球上的总摩擦力矩(不计外力.且假定全场压强是常数).

\noindent\textbf{8.13.} 设润滑油的密度 p=900kg / m" .运动粘性系数 $\nu$=O.37X10 Im\textasciitilde{}/s. 轴承

氏度 L=O.lm. 轴展周向速度 u = 12m/s. 人口间隙高度 h 1 =0. 0002m. 出口间隙高 度 h 2 =0. 000lm ，计算轴承承载力.

\section*{思考题}\addcontentsline{toc}{section}{思考题}

S8.1 在有端盖的圆柱筒内充满牛顿流体，分别讨论当该回柱作匀速直线运动 和绕轴线作匀速转动时.筒内流体作何运动?为什么?当容器为任意形状的封闭体 且作匀速直线运动时，容器内流体作何运动?为什么? S8.2 机翼分别在风洞和水洞中做实验，翼型的尺寸相同，要求翼型实验结果具 有相似性，实验中应满足什么条件? S8.3 在水池中做船舶阻力实验.若严格满足完全相似条件时，为什么要求模型 与原型的雷诺数和弗罗德数都相等? S8.4 模拟太宅中航天器的气动力实验，除了雷诺数和马赫数以外.是杳还需要 满足其他相似参数相等? S8.5 空气 If' 悬津颗粒的半径为 10 6 m. 在大气中运动时，是否可利用斯托克斯 近似计算作用力?
